\chapter{工程组织、交付与演进}
\label{chap:engineering}

\section{增量工程原则}

我们的实现不另建一套与 uCore 平行的“Agent 内核”，而是把身份、Context、工具、资源、观测和调度嵌入现有 PCB、syscall、VFS、IPC 与 scheduler 路径。业务策略保留在 Guest 用户态，provider 适配保留在 Host，而内核只处理有界二进制 ABI。这一原则使每个新功能都能回答：它改变了哪个既有系统责任，如何由内核强制，又如何被 Guest 消费。

\section{仓库组织}

\begin{table}[htbp]
  \centering
  \caption{主要代码目录与责任}
  \label{tab:repo-layout}
  \begin{tabularx}{\textwidth}{L{3.4cm}Y Y}
    \toprule
    目录 & 责任 & 交付证据 \\
    \midrule
    \code{os/agent_*.c} & 内核 identity、Context、工具、resource、metadata、observe、contract、fence & 参与生产内核对象链接，受 module/UAPI/stack gate 约束 \\
    \code{user/lib/} 与根目录、\code{user/include/} 中的 UAPI 头 & syscall wrapper、ABI 布局和公共常量 & kernel/user size/offset 静态断言 \\
    \code{user/src/} & 专项 Guest tests、\code{labdemo_ucore}、\code{agentlive_ucore}、\code{agentnexus_ucore} & 真实 RV64 应用与 QEMU markers \\
    \code{user/lib/agent_nexus.c} & N1 codec、role-filtered tool view、artifact 校验与 identity registry & 与 Nexus Guest 同步编译链接 \\
    \code{host_tools/} & CLI、daemon、provider relay、observer、strict validator & Host unit/mutation tests 和两窗口产品入口 \\
    \code{scripts/} & 边界 checker、stack check、QEMU runner & fail-closed 的静态和动态 gate \\
    \code{one_shot_metrics/} & 一次性采集计划、raw/tables/manifest、重绘和验证 & 30 boot、33 raw、19 tables、7,498 rows \\
    \code{docs/agentos/} & 架构、API、安全、验证、Console/Nexus 边界 & 实现与主张对应索引 \\
    \bottomrule
  \end{tabularx}
\end{table}

\section{四类产品入口}

四个入口各自回答不同问题，不能用模型的不确定性替代性能可复现性。

\begin{description}[leftmargin=3.1cm,style=nextline]
  \item[\code{make contest-demo}] 在同一 AgentOS Guest 中运行 traversal/compat 与 indexed/native 的固定 AB/BA 工作流，是可比较的竞赛性能基线。
  \item[\code{make dual-platform-run}] 构建并运行 Plain uCore 与 AgentOS-uCore 双目标，用于验证跨目标业务合同和工作量对齐。
  \item[\code{make agent-live-demo}] 运行单 Agent 的 Guest model/tool loop；默认使用 replay，live provider 需显式选择。
  \item[\code{make agentos-nexus}] 以 live provider 启动 daemon 与 controller；只读窗口由 \code{make agentos-nexus-observe} 单独启动。
\end{description}

对 Nexus，\code{agentos-nexus-image} 只构建镜像，\code{agentos-nexus-check} 只运行本地静态/Host 合同，\code{agentos-nexus-replay} 启动真实 QEMU strict replay。\code{agentos-nexus} 默认选择 DeepSeek/live provider，\code{agentos-nexus-deepseek} 是其显式别名；两者都依赖人工提供的网络与 API key，且均需另启 observer 窗口。这些入口名称同时约束证据层次。

\section{构建与回归}

\subsection{建议顺序}

\begin{lstlisting}[language=bash,caption={从快速合同到 QEMU 产品的分层验证}]
# Host / protocol / static contracts
make local-host-selftests
make agent-uapi-check TOOLPREFIX=riscv64-linux-gnu-
make agent-module-check TOOLPREFIX=riscv64-linux-gnu-

# Real production objects and stack budgets
make build TOOLPREFIX=riscv64-linux-gnu-
make kernel-stack-check TOOLPREFIX=riscv64-linux-gnu-
make user-stack-check TOOLPREFIX=riscv64-linux-gnu-

# Product-specific contracts and QEMU replay
make agentos-nexus-check
make agentos-nexus-image TOOLPREFIX=riscv64-linux-gnu-
make agentos-nexus-replay TOOLPREFIX=riscv64-linux-gnu-
\end{lstlisting}

不同本地 RISC-V 工具链可能使用 \code{riscv-none-elf-} 或 \code{riscv64-linux-gnu-}，实际以当前仓库 Makefile 和环境为准。Windows 下的现场产品通常通过 WSL 运行 make/QEMU，Host CLI 与 key 路径依照已有入口转换；不将 API key 传入 Guest 镜像。

\subsection{共享输出约束}

Console/Nexus image 复用 \code{build/kernel} 与 \code{nfs/fs*.img}，因此不应在同一 worktree 并发构建两个镜像。性能采集还必须避免与 observer 或其他 QEMU 实例并发，以免改变 Host 负载和结果口径。一次性数据目录已经写入 \code{COMPLETED}，日常 CI 不会重跑或覆盖它。

\section{交付物和可追溯性}

竞赛交付不只包含一个可启动镜像，还包含：

\begin{itemize}
  \item 可从源码重建的 Plain/AgentOS/Nexus Guest 目标；
  \item UAPI 头文件、生产对象清单与调用图栈 gate；
  \item 确定性 contest-demo 与 strict Nexus replay fixture；
  \item Host controller/observer 事件、Guest 串口输出和 validator 结果；
  \item one-shot raw、manifest、long tables、validation warning 和可重绘图表；
  \item 架构、API、安全、验证、Nexus 操作和本决赛报告。
\end{itemize}

关键主张能够通过“文档章节 $\rightarrow$ 实现文件 $\rightarrow$ checker/QEMU 入口 $\rightarrow$ raw/table/fixture”回溯。附录\ref{app:source-index}给出高信号源码索引，附录\ref{app:reproduction}给出复验步骤与数据解释。

\section{从确定性流程到 Nexus 的工程演进}

第一阶段以 \code{labdemo_ucore} 建立 Orchestrator/Sentinel/Investigator/Recovery 的真实多 Agent 身份、route、Context 和文件工作流，为 traversal/indexed 配对比较提供可重复基线；Plain uCore 与 AgentOS-uCore 的双目标对齐由 \code{rp_*} Guest 和 \code{dual-platform-run} 另行验证。第二阶段的 \code{agentlive_ucore} 建立 model request/response、真实 typed V2 工具和结果回灌的闭环。第三阶段将该循环扩展为长驻 CLI，再用 Nexus 统一动态委派、artifact 和独立 observer。

这条演进线保留了旧路径的用途：\code{contest-demo} 继续是性能基准，不被模型不确定性替代；\code{agent-live-demo} 仍是更小的 model-loop 验收；Nexus 才是现场交互主演示。

\section{后续演进路径}

下列方向是建立在当前边界上的工程计划，不是已完成能力：

\begin{enumerate}
  \item 为 artifact store 增加内核或可验签名的 object ACL/seal，并设计跨 boot 恢复与容量回收；
  \item 让 Task Channel 承载有 generation-safe handle 的真实业务 input/result，保留 cancellation 与 exactly-one completion；
  \item 将 Guest 调用级审批升级为内核可验证的有界 token，同时保留人机交互体验；
  \item 为外部资源增加 Host gateway 工具，将网络输出统一标记为不可信，不让 Host 获得业务决策权；
  \item 在保持 Nexus N1 和内核 ABI 分层的前提下，实现真实 MCP stdio/HTTP Client/Server 互操作；
  \item 在多 Hart 与更大工作流上扩展 EEVDF、观测和资源账户实验，增加统计独立 boot 而不复制 inner pair。
\end{enumerate}

每一步都应先定义责任边界与可执行验收，再增加 API。我们不为了展示丰富度而增加没有真实 consumer 的内核接口。
